\section{Question: the same object --- what HLV does with it stays HLV's business}
One statement stands at the outset and is meant to be unambiguous: FFGFT and HLV work with the
\emph{same} geometric object. This is not an interpretation and not a thesis, but a computable
fact, established below --- the $\varphi$-icosahedral axis system on which $\theta=2/9$ sits as an
invariant. Everything else in this document is a consequence of this single finding.

Strictly to be separated from it is a second question, which is here \emph{not} asked and
\emph{not} answered: what HLV does with this object --- which vortex or edge-mode readings, which
derivations and which formulations HLV builds upon it. That is HLV's business, not FFGFT's.
Whether those formulations are correct is here neither claimed nor tested nor addressed. FFGFT
establishes the shared object and --- via $E_0$ --- provides a test direction; for the use HLV
makes of the object, HLV alone is responsible.

The purpose of the document is thereby deliberately two-fold and modest. It records that both
compute over the same object, and that FFGFT offers a way to measure HLV against the empirical
data; and it draws at the same time the boundary beyond which nothing is asserted about HLV.

\section{The shared $\varphi$-icosahedral object}
In the available computational representation of HLV --- the cut-and-project of the homology
check script \texttt{ffgft\_hlv\_homology\_check\_v3.py} --- physical space is generated by the
icosahedral projector with parameter $\tau=\varphi$, $\varphi=\tfrac{1+\sqrt5}{2}$. Its six
projected axes are the six five-fold axes of the icosahedron: all fifteen pairwise angles equal
$\arccos(1/\sqrt5)=63{,}435^\circ$. On exactly these axes lives FFGFT's $C_3$-in-$A_5$ doubling
(Dok.~285, Dok.~293).

It is therefore not two similar structures but the same object. Applying to these axes the
five-fold redistribution of the $C_3$ modes (Dok.~293), the share into the trivial mode is
exactly $p_0=2/9$. The Koide angle $\theta=2/9$ is thereby an \emph{invariant on the shared
object}, not a translation of one framework into the other. The value is moreover specific to the
golden ratio: only $\tau=\varphi$ makes the angle spectrum exactly uniform and the trivial share
exactly $2/9$; the silver value $\tau=1+\sqrt2$ gives $0{,}2534$ and an already measurably shifted
spectrum (angles $45^\circ$ and $69{,}3^\circ$). The shared object is thus not just any
cut-and-project, but the $\varphi$-icosahedral one.

This statement is computational and free of circularity: it follows from the comparison of the
projector axes and the mode redistribution, checkable in \texttt{ikosaeder\_theta\_delta.py}
(seed $20780458$) and in the reproducible angle-spectrum check
\texttt{a5\_bridge\_discrimination.py}.

\section{Delimitation: which HLV object is shared}
The shared object established here is the \emph{geometric} layer of HLV --- the
$\varphi$-icosahedral axis system on which $\theta=2/9$ sits as an invariant. Strictly to be
separated from it is a second HLV object of the same name: the \emph{temporal-dynamic} spiral
time, a non-Markov operator with triadic time coordinate $\Psi(t)=(t,\phi(t),\chi(t))$ and a
memory channel (Zenodo 20643462). This temporal-dynamic spiral time carries neither $A_5$ nor the
icosahedral projector; it is \emph{not} the object shared here and does not fall under the present
finding.

The test direction of this document concerns the geometric layer alone. The comparison of the
temporal-dynamic spiral time with FFGFT's time-winding is carried out separately (Dok.~295); there
a computational check shows that this object does not carry the self-similarity ratio $e$ of the
FFGFT log spiral --- on its own definitions it is a winding of constant pitch, not equiangular.
Both delimitations together keep the talk of \enquote{shared} clean: what is shared is the
$\varphi$-icosahedral geometric object, not HLV's spiral time as a time operator. The shared
invariant between the two layers is the $Z_3$ alone ($C_3<A_5$); the dimension counts of the two
frameworks ($3$, $6$) coincide, their meaning does not (Dok.~295).

\section{The $E_0$ asymmetry}
The shared object is scale-free. The geometry knows neither mass nor energy; it yields pure
numbers --- the angle and the denominator-$9$ weights (Dok.~293). Here the actual asymmetry
between the two frameworks sets in.

FFGFT possesses a reference point $E_0$ (Dok.~292) that binds the scale-free pattern to the
measured leptons $e,\mu,\tau$ and fixes the absolute position. FFGFT can thus \emph{compute back
empirically} from the pattern: from the shared object, the denominator-$9$ weights and $E_0$,
absolute masses follow. HLV possesses no such scale. It carries the pattern --- the shared
geometric object --- but no anchor in energy or mass. HLV can therefore obtain no masses from the
object alone. It can only do so if it adopts $E_0$ from FFGFT; in that case, however, it is
FFGFT's empirical content that computes, merely expressed in HLV's coordinates.

\section{The test direction: from FFGFT to HLV}
From the shared geometry and the one-sided scale a directed relation follows. Because FFGFT
supplies the missing scale, FFGFT offers a way to test the geometric content of HLV against the
data: one equips the shared object with $E_0$ and tests whether it reproduces the lepton masses.
This is a verification direction from FFGFT to HLV.

It is expressly \emph{not} a confirmation that HLV's own formulations are correct. The shared
object states only that both use the same geometry --- not that the interpretations and
derivations HLV erects on this geometry are valid. The test direction tests the correspondence of
geometry and masses with FFGFT's means; it does not touch HLV's own construction and does not
ennoble it.

Conversely: HLV supplies FFGFT empirically with nothing it does not already have. The pattern is
shared, the scale FFGFT has itself. There is no empirical transfer from HLV to FFGFT, and FFGFT
does not need HLV for the masses. The empirical arrow is one-sided.

\section{What the finding says and what it does not}
Established and computationally secured is: the shared $\varphi$-icosahedral object; $\theta=2/9$
as its invariant; the specificity for $\tau=\varphi$; and the $E_0$ asymmetry, from which the
test direction follows.

Not established and not claimed is: the correctness of the formulations HLV contains; an
independent empirical confirmation of FFGFT by HLV; and the question of whether the icosahedral
projector is \emph{intrinsic} to HLV's fundamental claim or constitutes a modelling choice in the
check scripts. The latter is for HLV itself to clarify and is here left open.

\section{Methodological placement}
Within FFGFT's three methodological layers the finding places clearly. The shared object is a
bridge finding: a geometric, algebraically demonstrable correspondence. The back-computation to
masses is a core derivation, it follows from $\xi$ and $E_0$ and is FFGFT-internal. HLV's
contribution to the empirical content is zero. In the sense of P35 the statement that HLV is
correct remains a candidate and not a derivation; derived is the shared object alone.

\section{The first check attempt is superseded}
The first check attempt --- the homology check \texttt{ffgft\_hlv\_homology\_check\_v3.py}, a
coarse topological point-cloud diagnostic --- is superseded by the exact Hilbert-space variant
(Dok.~293) and was in any case structurally blind to the load-bearing question: the persistent
$H_1$ of a point cloud does not see the axis-angle structure on which $2/9$ hangs, so that
$\tau=\varphi$ (exactly $2/9$) and the control value $\tau=\sqrt2$ ($0{,}2139$) produce
topologically nearly indistinguishable clouds. Its earlier finding of \enquote{only weak
separation from generic} is therefore not a statement about $\varphi$ or $A_5$, but the
expectable blindness of the coarse statistic. In the spectral variant, by contrast, $\theta=2/9$
appears directly as the exact invariant $p_0=\lvert\langle v_0|R_5 v_e\rangle\rvert^2=2/9$; the
shared-object finding rests on it, not on the homology diagnostic.

\section{The second test: what it resolves and where it becomes circular}
The second check attempt --- the instrument \texttt{a5\_bridge\_discrimination.py} (seed
$20780458$) --- uses exactly the statistic the first lacked: the angle spectrum of the six
projected axes, measured as distance to the icosahedral reference spectrum in which all fifteen
pairwise angles equal $63{,}435^\circ$. It thereby resolves the axis-angle structure on which
$2/9$ hangs. It is calibrated so that an exactly icosahedral object gives distance zero, while
generic random tight frames and a truly generic $\sqrt2$ construction lie in a broad generic
range. It is thus the instrument that can see $\varphi$ and $A_5$ where the first was blind.

What the second test yields on the available representation is two-fold. First, the HLV projector
with $\tau=\varphi$ reads exactly icosahedral: all fifteen angles $63{,}435^\circ$, distance zero,
and the five-fold redistribution on its axes gives exactly $2/9$. This, however, is a consequence
of the construction --- the object \emph{is} the icosahedral projector --- so the reading is
circular and no independent confirmation of $A_5$. The instrument correctly treats this not as a
verdict but makes the circularity explicit. Second, it demonstrates quantitatively that the
earlier \enquote{generic} control with $\tau=\sqrt2$ is not generic: RMS distance $3{,}47$ to the
icosahedral reference spectrum (the exactly icosahedral object has distance $0$), a two-valued
spectrum $61{,}87^\circ/70{,}53^\circ$, and a redistribution weight $0{,}2139$ that is not $2/9$.
This distance lies far below the generic band of random frames (median RMS about $22$,
$5$--$95\,\%$ roughly $17$--$28$): the $\sqrt2$ control is thus structured and not generic. The
$\varphi$-specificity is thereby nailed down: only $\tau=\varphi$ makes the spectrum exactly
uniform and the weight exactly $2/9$.

The yield of the second test is thus diagnostic and instrumental, not confirming. It provides the
correct, reproducible statistic that resolves what the first could not; it quantifies the
circularity of a check of the available object and the contamination of the earlier control; and
it isolates $\varphi$ as the distinguished value. What it does not yet provide is an independent
confirmation: that would require an HLV object represented \emph{without} a baked-in icosahedral
projector. Until such an object is available, the instrument stands ready, and the shared-object
finding rests on the exact spectral invariant, not on a non-circular discrimination.

\section{Status}
Secured is the shared $\varphi$-icosahedral object and the $E_0$ asymmetry, and from it a test
direction from FFGFT to HLV: FFGFT offers the possibility of testing the geometric content of HLV
against the measured masses. No correctness of the formulations HLV contains is claimed. One point
remains open: whether the icosahedral projector is intrinsic to HLV. Bound to it is the only still
possible independent check --- a non-circular discrimination requiring an HLV object without a
baked-in icosahedral projector; the instrument \texttt{a5\_bridge\_discrimination.py} stands ready
for it. The earlier homology contrast is settled, not open.

$\theta=2/9$ itself stands independent of all this. It follows from the $Z_3$ circulant
diagonalisation and the icosahedral redistribution (Dok.~293) and needs HLV in no way. The
relation recorded here is a shared geometry with a one-sided test direction, not a mutual
confirmation.
